% Source: https://artofproblemsolving.com/community/c3h15248_mathcounts_contest_1_sprint_questions
% Stable identity: aops-topic-15248-q06
Let $O$ be the total amount of orange juice; then there are $80-O$ ounces of cranberry juice, and the ratio $\dfrac{80-O}{O}=\dfrac{80}{O}-1$ is largest when $O$ is smallest.

With two orange bottles, one holds $16$ oz and the other at least $12$ oz, so $O\ge28$; with three orange bottles $O\ge16+12+12=40$. So $O\ge28$.

$O=28$ is possible: orange bottles of $16$ and $12$ oz, and three cranberry bottles holding $52$ oz in total, for example $20,20,12$ oz (each between $12$ and $20$). The largest ratio is
\[\dfrac{52}{28}=\boxed{\dfrac{13}{7}}.\]
